% Part: computability
% Chapter: computability-theory
% Section: fixed-point-thm

\documentclass[../../../include/open-logic-section]{subfiles}

\begin{document}

\olfileid{cmp}{thy}{fix}
\olsection{স্থির-বিন্দু উপপাদ্য}

থামার সমস্যাটি আবার বিবেচনা করা যাক। সাময়িক সংকেত হিসেবে
$\tuple{x, y}$-এর বদলে $\gn{\cfind{x}(y)}$ লিখি; এটিকে
$\cfind{x}(y)$ মানটির একটি “নাম”-এর প্রতিনিধিত্ব হিসেবে ভাবো। এই সংকেতে
থামার সমস্যার অনির্ণেয়তা নিয়ে আমাদের একটি প্রমাণকে নতুনভাবে বলা যায়।

প্রশ্ন: নিচের বৈশিষ্ট্যসহ কোনো গণনসাধ্য অপেক্ষক~$h$ আছে কি? প্রত্যেক
$x$ ও $y$-এর জন্য
\[
h(\gn{\cfind{x}(y)}) =
\begin{cases}
1 & \text{যদি $\cfind{x}(y) \fdefined$} \\
0 & \text{অন্যথায়।}
\end{cases}
\]

উত্তর: না; তা হলে আংশিক অপেক্ষক
\[
g(x) \simeq
\begin{cases}
0 & \text{যদি $h(\gn{\cfind{x}(x)}) = 0$} \\
\text{অসংজ্ঞায়িত} & \text{অন্যথায়}
\end{cases}
\]
গণনসাধ্য হতো, ফলে তার কোনো সূচক~$e$ থাকত। কিন্তু তখন
\[
\cfind{e}(e) \simeq
\begin{cases}
0 & \text{যদি $h(\gn{\cfind{e}(e)}) = 0$} \\
\text{অসংজ্ঞায়িত} & \text{অন্যথায়,}
\end{cases}
\]
অর্থাৎ $\cfind{e}(e)$ সংজ্ঞায়িত যদি এবং কেবল যদি সেটি সংজ্ঞায়িত নয়;
এটি একটি বিরোধ।

এবার $\cfind{e}$-সহ সমীকরণটি দেখো। এক অর্থে সেখানে স্ব-উল্লেখ আছে:
$\cfind{e}(e)$-এর মানটি যেন একটি নির্দিষ্টভাবে তার নিজের নাম
$\gn{\cfind{e}(e)}$-এর উপর নির্ভর করে, আমরা সেই ব্যবস্থা করেছি।
স্থির-বিন্দু উপপাদ্য বলে, সাধারণভাবেই আমরা এটি \emph{করতে পারি}—শুধু
বিরোধ প্রমাণ করার জন্য নয়।

\olref{lem:fixed-equiv} স্থির-বিন্দু উপপাদ্য বলার দুটি সমতুল্য উপায় দেয়।
যুক্তিগতভাবে, বক্তব্য দুটি সমতুল্য হওয়া তাদের উভয়ের সত্যতা থেকেই অনুসৃত;
কিন্তু আমাদের আসল অর্থ হলো, প্রতিটি বক্তব্য অন্যটি থেকে সরাসরি অনুসরণ
করে, তাই একই উপপাদ্যের বিকল্প বক্তব্য হিসেবে তাদের নেওয়া যায়।

\begin{lem}
\ollabel{lem:fixed-equiv}
নিচের বক্তব্যগুলি সমতুল্য:
\begin{enumerate}
\item প্রত্যেক আংশিক গণনসাধ্য অপেক্ষক $g(x,y)$-এর জন্য এমন একটি সূচক~$e$
  আছে, যাতে প্রত্যেক~$y$-এর জন্য
\[
\cfind{e}(y) \simeq g(e,y).
\]
\item প্রত্যেক গণনসাধ্য অপেক্ষক~$f(x)$-এর জন্য এমন একটি সূচক~$e$ আছে,
  যাতে প্রত্যেক~$y$-এর জন্য
\[
\cfind{e}(y) \simeq \cfind{f(e)}(y).
\]
\end{enumerate}
\end{lem}

\begin{proof}
$(1) \Rightarrow (2)$: $f$ দেওয়া থাকলে $g(x,y) \simeq
\fn{Un}(f(x),y)$ দিয়ে~$g$-কে সংজ্ঞায়িত করো। (1) ব্যবহার করে এমন একটি
সূচক~$e$ নাও, যাতে প্রত্যেক~$y$-এর জন্য
\begin{align*}
\cfind{e}(y) & \simeq \fn{Un}(f(e),y) \\
& \simeq \cfind{f(e)}(y).
\end{align*}

$(2) \Rightarrow (1)$: $g$ দেওয়া থাকলে $s$-$m$-$n$ উপপাদ্য ব্যবহার
করে এমন~$f$ নাও, যাতে প্রত্যেক $x$ ও~$y$-এর জন্য
$\cfind{f(x)}(y) \simeq g(x,y)$। (2) ব্যবহার করে এমন একটি সূচক~$e$ নাও,
যাতে
\begin{align*}
\cfind{e}(y) & \simeq \cfind{f(e)}(y) \\
& \simeq g(e,y).
\end{align*}
এতে প্রমাণ শেষ হলো।
% BN-SRC-158 TARGET-CORRECTION-BEGIN
\par\smallskip\noindent\textnormal{\small\textbf{উৎস-সংশোধন (BN-SRC-158):} হিমায়িত ইংরেজি সমতুল্যতা-প্রমাণের দুটি align ব্লকে আংশিক অপেক্ষকের চারটি সমতার জন্য সাধারণ “$=$” ব্যবহার করেছে। বিবৃতিতেই থাকা যুগপৎ সংজ্ঞায়িততার সমতা অনুসারে বাংলা পাঠে চারটিতেই “$\simeq$” রাখা হয়েছে।} % BN-SRC-158 TARGET-CORRECTION-END
\end{proof}

\begin{explain}
(1) সত্য (এবং তাই (2)-ও সত্য) দেখানোর আগে বক্তব্যটি কত অদ্ভুত তা ভাবো।
$e$-কে একটি কম্পিউটার প্রোগ্রাম হিসেবে ভাবো; (1) বলে, যেকোনো আংশিক
গণনসাধ্য $g(x,y)$ দেওয়া থাকলে এমন একটি কম্পিউটার প্রোগ্রাম~$e$ পাওয়া
যায়, যা $g_e(y) \simeq g(e,y)$ গণনা করে। অন্যভাবে বললে, এমন একটি
কম্পিউটার প্রোগ্রাম পাওয়া যায়, যা প্রোগ্রামটিকেই উল্লেখ করে এমন একটি
অপেক্ষক গণনা করে।
\end{explain}

\begin{thm}
\olref{lem:fixed-equiv}-এর দুটি বক্তব্যই সত্য। বিশেষ করে, প্রত্যেক আংশিক
গণনসাধ্য অপেক্ষক $g(x,y)$-এর জন্য এমন একটি সূচক~$e$ আছে, যাতে
প্রত্যেক~$y$-এর জন্য
\[
\cfind{e}(y) \simeq g(e,y).
\]
\end{thm}

\begin{proof}
উপরের থামার সমস্যার আলোচনায় উপাদানগুলি ইতিমধ্যেই নিহিত আছে।
$\fn{diag}(x)$-কে এমন একটি গণনসাধ্য অপেক্ষক ধরি, যা প্রত্যেক~$x$-এর
জন্য $f_x(y) \simeq \cfind{x}(x,y)$ অপেক্ষকটির একটি সূচক ফেরত দেয়,
অর্থাৎ
\[
\cfind{\fn{diag}(x)}(y) \simeq \cfind{x}(x,y).
\]
$\fn{diag}$-কে এমন একটি অপেক্ষক হিসেবে ভাবো, যা একটি দ্বিস্থানী অপেক্ষকের
প্রোগ্রামকে একস্থানী অপেক্ষকের প্রোগ্রামে রূপান্তর করে; মূল প্রোগ্রামটিকেই
তার প্রথম আর্গুমেন্ট হিসেবে স্থির করে নতুন প্রোগ্রামটি পাওয়া যায়।
$\fn{diag}$-কে আনুষ্ঠানিকভাবে এভাবে সংজ্ঞায়িত করা যায়: প্রথমে~$s$-কে
সংজ্ঞায়িত করি,
\[
s(x,y) \simeq \fn{Un}^2(x,x,y),
\]
যেখানে $\fn{Un}^2$ হলো আংশিক গণনসাধ্য দ্বিস্থানী অপেক্ষকগুলির জন্য
সার্বজনীন একটি ত্রিস্থানী অপেক্ষক। এরপর $s$-$m$-$n$ উপপাদ্য অনুসারে এমন
একটি আদিম পুনরাবৃত্ত অপেক্ষক~$\fn{diag}$ পাওয়া যায়, যাতে
\[
\cfind{\fn{diag}(x)}(y) \simeq s(x,y).
\]

এবার অপেক্ষক~$l$-কে সংজ্ঞায়িত করি:
\[
l(x,y) \simeq g(\fn{diag}(x),y).
\]
এবং $\gn{l}$-কে~$l$-এর একটি সূচক ধরি। শেষে
$e = \fn{diag}(\gn{l})$ ধরি। তাহলে প্রত্যেক~$y$-এর জন্য
\begin{align*}
\cfind{e}(y) & \simeq \cfind{\fn{diag}(\gn{l})}(y) \\
& \simeq \cfind{\gn{l}}(\gn{l}, y) \\
& \simeq l(\gn{l}, y) \\
& \simeq g(\fn{diag}(\gn{l}),y) \\
& \simeq g(e, y),
\end{align*}
যা প্রয়োজন ছিল।
\end{proof}

\begin{explain}
কী ঘটছে? ধরা যাক, তোমাকে নিজেকেই মুদ্রণ করে এমন একটি কম্পিউটার প্রোগ্রাম
লিখতে বলা হয়েছে। আরও ধরা যাক, তুমি এমন একটি প্রোগ্রামিং ভাষায় কাজ করছ,
যার প্রতীকক্রম-অপেক্ষকের সমৃদ্ধ ও বিচিত্র সংগ্রহ আছে। বিশেষ করে, ধরা যাক
তোমার প্রোগ্রামিং ভাষায় নিচের মতো কাজ করে এমন একটি~$\fn{diag}$ অপেক্ষক
আছে: নিবেশ হিসেবে প্রতীকক্রম~$s$ দেওয়া হলে $\fn{diag}$, $s$-এ থাকা
`x' প্রতীকের প্রতিটি দৃষ্টান্ত খুঁজে মূল প্রতীকক্রমটির উদ্ধৃত সংস্করণ দিয়ে
সেটিকে প্রতিস্থাপন করে। যেমন, নিবেশ হিসেবে প্রতীকক্রম
\begin{quote}
\begin{verbatim}
hello x world
\end{verbatim}
\end{quote}
দিলে অপেক্ষকটি নির্গম হিসেবে ফেরত দেয়
\begin{quote}
\begin{verbatim}
hello 'hello x world' world
\end{verbatim}
\end{quote}
। সে ক্ষেত্রে কাঙ্ক্ষিত প্রোগ্রামটি লেখা সহজ; যাচাই করা যায়,
\begin{quote}
\begin{verbatim}
print(diag('print(diag(x))'))
\end{verbatim}
\end{quote}
কাজটি করে। C++ ও Java-র মতো বেশি প্রচলিত প্রোগ্রামিং ভাষাতেও একই ধারণা
(আরও জটিল বাস্তবায়নসহ) কাজ করে।

স্থির-বিন্দু উপপাদ্যের প্রমাণ থেকে আমরা মাত্র কয়েক ধাপ দূরে। ধরা যাক,
মুদ্রণ অপেক্ষকের একটি রূপ~$\fn{print}(x,y)$ একটি প্রতীকক্রম~$x$ এবং আরেকটি
সাংখ্যিক আর্গুমেন্ট~$y$ নেয়, তারপর প্রতীকক্রম~$x$-কে $y$ বার মুদ্রণ করে।
তাহলে “প্রোগ্রাম”
\begin{quote}
\begin{verbatim}
getinput(y);
print(diag('getinput(y); print(diag(x), y)'), y)
\end{verbatim}
\end{quote}
নিবেশ~$y$-এ নিজেকেই $y$ বার মুদ্রণ করে।
$\fn{getinput}$---$\fn{print}$---$\fn{diag}$ কাঠামোটিকে যেকোনো অপেক্ষক
$g(x,y)$ দিয়ে প্রতিস্থাপন করলে পাওয়া যায়
\begin{quote}
\begin{verbatim}
g(diag('g(diag(x), y)'), y)
\end{verbatim}
\end{quote}
যা এমন একটি প্রোগ্রাম, যে নিবেশ~$y$-এ প্রোগ্রামটি নিজে এবং~$y$-কে নিয়ে
$g$ চালায়। “উদ্ধৃত করা”-কে “একটি সূচক ব্যবহার করা” হিসেবে ভাবলে উপরের
প্রমাণটি পাওয়া যায়।

এখনকার মতো প্রমাণটিকে আনুষ্ঠানিক কারসাজি বা কালো জাদু ভাবতে চাইলে ক্ষতি
নেই। তবে উপরে দেওয়া যুক্তির বিস্তারিত পুনর্গঠন করতে পারা উচিত। আমরা যখন
অসম্পূর্ণতা উপপাদ্যগুলি (এবং সংশ্লিষ্ট “স্থির-বিন্দু উপপাদ্য”) প্রমাণ করব,
তখন এটি কেন কাজ করে তা বোঝার অন্য উপায় আলোচনা করব।
\end{explain}

\begin{tagblock}{lambda}
\begin{digress}
একই ধারণা ব্যবহার করে একটি “স্থির-বিন্দু” সমাবেশক পাওয়া যায়। ধরা যাক,
তোমার কাছে একটি ল্যাম্বডা পদ~$g$ আছে এবং তুমি এমন আরেকটি পদ~$k$ চাও,
যার বৈশিষ্ট্য হলো $k$, $gk$-এর সঙ্গে $\beta$-সমতুল্য। আমাদের সাংকেতিক
রীতি ব্যবহার করে পদগুলি সংজ্ঞায়িত করি:
\[
\fn{diag}(x) = xx
\]
এবং
\[
l(x) = g(\fn{diag}(x))
\]
অন্যভাবে বললে, $l$ হলো $\lambd[x][g(xx)]$ পদটি। $k$-কে $ll$ পদটি ধরি।
তাহলে
\begin{align*}
k & = (\lambd[x][g(xx)])(\lambd[x][g(xx)]) \\
& \red  g((\lambd[x][g(xx)])(\lambd[x][g(xx)])) \\
& = gk.
\end{align*}
যদি
\[
Y = \lambd[g][((\lambd[x][g(xx)])(\lambd[x][g(xx)]))]
\]
নেওয়া হয়, তবে $Yg$ ও $g(Yg)$ একটি সাধারণ পদে হ্রাস পায়; তাই
$Yg \equiv_\beta g(Yg)$। এটি “কারির সমাবেশক” নামে পরিচিত। তার বদলে যদি
\[
Y = (\lambd[xg][g(xxg)])(\lambd[xg][g(xxg)])
\]
নেওয়া হয়, তবে আসলে $Yg$, $g(Yg)$-তে হ্রাস পায়, যা আরও শক্তিশালী
বক্তব্য। $Y$-এর এই শেষ সংস্করণটি “টুরিংয়ের সমাবেশক” নামে পরিচিত।
\end{digress}
\end{tagblock}

\end{document}
